\section{Related Work}
\label{related_work}

This section positions the study within requirement elicitation, context-aware requirements engineering, and human-centered intelligent requirement analysis.

\subsection{Traditional Requirement Elicitation}
\label{rw:traditional_elicitation}

This subsection reviews interviews, questionnaires, workshops, scenarios, and prototyping as established elicitation techniques.
The discussion should emphasize their reliance on stakeholder communication and the difficulty of systematically capturing implicit environmental assumptions.
Prior work on requirement specification further shows that reducing ambiguity while preserving stakeholder readability remains a persistent challenge~\cite{darifControlledNaturalLanguage2026}.

\subsection{Context-Aware Requirement Engineering}
\label{rw:context_aware_re}

This subsection covers approaches that model user context, operational context, and situational variability in requirements.
The comparison should clarify how this paper targets context discovery during elicitation rather than downstream modeling.

\subsection{Environment-Aware Software Systems}
\label{rw:environment_aware_systems}

This subsection reviews software systems that adapt to deployment environments, devices, sensors, connectivity, or runtime conditions.
The relevance is to identify what environmental information should be elicited before implementation decisions are made.

\subsection{LLM-Based Requirement Elicitation}
\label{rw:llm_elicitation}

This subsection discusses recent work that uses language models to assist requirement generation, refinement, summarization, or questioning.
The gap to emphasize is that fluent generated requirements may still miss deployment context and operational constraints.

\subsection{Multi-Agent Collaboration}
\label{rw:multi_agent}

This subsection reviews collaborative agent settings in which multiple roles decompose analysis tasks.
The discussion should stay focused on role separation for elicitation and context analysis, not on complex agent architectures.

\subsection{Human-in-the-Loop Requirement Engineering}
\label{rw:hitl}

This subsection reviews stakeholder-involved requirement analysis and clarification.
The paper positions human-in-the-loop interaction as a mechanism for validating uncertain environmental assumptions and completing missing context.
